\paragraph{Vue.js}
Vue.js \cite{wikipedia_vue} es un framework desarrollado en JavaScript. Su arquitectura está basado en componentes, lo que 
sginifica que cada elemento de la interfaz encapsula su lógica propia y estilo permitiendo una reutilización 
del código en distintas partes de la aplicación. La curva de aprendizaje es baja, lo cual facilita el
desarrollo para aplicaciones sencillas.Además, cuenta con herramientas integradas que simplifican la interacción 
con formularios y el manejo del estado.

\paragraph{React}   
React \cite{reactwiki2015} es una biblioteca de Javascript basado en componentes. Utiliza JSX, una sintaxis que combina HTML con Javascript 
facilitando la creación de componentes reutilizados. Cuenta un alto rendimiento al actualizar elementos de la interfaz.
A diferencia del resto de frameworks, React se enfoca exclusivamente en la 
interfaz del usuario, por lo que requiere incorporar librería adicionales para funcionalidades como
el enrutamiento o la gestión de formularios. 

\paragraph{Angular}
Angular \cite{wikipedia_angular} es un framework desarrollado en TypeScript basado en componentes. A diferencia de Vue.js y React, Angular destaca por una estructura 
más rígida al proyecto, lo que facilita el mantenimiento en aplicaciones grandes. El uso de TypeScript, proporciona
un tipado estático, por lo que mejora la detección de errores durante el desarrollo. Además, ofrece un sistema 
de comunicación entre componentes y dispone de un sistema de inyección de dependencias. Sin embargo, su curva de aprendizaje
es más alta respecto a las otras dos.

Tras comparar los frameworks frontend, se ha optado por Angular debido a su tipado estático proporcionado por TypeScript y su robusto 
sistema de comunicación e inyección de dependencias entre componentes, características que garantizan mayor fiabilidad y escalabilidad en el desarrollo de la aplicación.

